\begin{abstract}
A multi-task leaderboard turns per-task scores into one ranking, and a developer can try to move a chosen model to first place by training on data related to some of the tasks. We ask how many score points, over how many tasks, first place requires, and how this depends on the aggregation rule. Investment is the points raised on chosen tasks up to a per-task cap, at a linear, task-symmetric cost; we compute its exact minimum for the mean, geometric, harmonic and trimmed means, median, minimum, Borda, plurality and the Condorcet-winner requirement on HELM MMLU (57 subjects, 79 models) and the Open LLM Leaderboard (6 tasks, 4497 models). Two findings hold on both datasets: the minimum rule is cheaper than the mean in points for every target, most for balanced targets (the cost ratio rises with the coefficient of variation of the target's scores), and rank alone explains most of the variance in cost, the gap to the best rival adding explanatory power only for some rules. Other findings are conditional. For targets ranked 11--30 the median uniform per-task gain needed under the mean is \hstarBinHelmMeanElevenThirty{} points (HELM) and \hstarBinOllbMeanElevenThirty{} points (Open LLM Leaderboard), against a pooled 95th-percentile fine-tuning gain of \empPooled{} points, whereas the median over all Open LLM Leaderboard targets, \ollbShareOverHundred\% of which rank beyond 100, is \hstarOllbMean{}; under per-task calibrated caps most targets cannot reach first place, and for HELM this depends on how the gains are transferred (\empMmluQ{} or \empPooled{} points); and the order of rules under fixed gains differs between datasets and, on the Open LLM Leaderboard, between data slices. We mark derived identities apart from measurements and report what failed or was not done.
\end{abstract}
